\documentclass[11pt]{letter}

\usepackage[T1]{fontenc}
\usepackage[utf8]{inputenc}
\usepackage[margin=1in]{geometry}
\usepackage{hyperref}
\hypersetup{colorlinks=true, urlcolor=blue, linkcolor=blue}

\signature{Quang-Vinh Dang\\Corresponding author\\British University Vietnam, Hung Yen, Vietnam\\vinh.dq4@buv.edu.vn\\[2ex]
Thi-Hong-Hanh Nguyen\\Banking Academy Vietnam, Hanoi, Vietnam}
\address{Quang-Vinh Dang \\ British University Vietnam \\ Hung Yen, Vietnam \\ vinh.dq4@buv.edu.vn}
\date{\today}

\begin{document}

\begin{letter}{Dr.\ Ali Kutan, Editor\\\textit{Borsa Istanbul Review}}

\opening{Dear Dr.\ Kutan,}

We submit ``Credible on Paper? Green Banking Policy, Access to Finance, and Institutional Capacity in the Greening of Firms Worldwide'' for consideration as an original research paper in \textit{Borsa Istanbul Review}.

The paper asks whether green banking policy adoption --- specifically, membership of the IFC-convened Sustainable Banking and Finance Network, which now spans more than sixty emerging and developing economies --- changes the firm-level relationship between access to bank credit and green investment behavior, and whether that depends on the regulatory capacity available to make the supervisory guideline credible. It is, to our knowledge, the first firm-level test of that premise across a broad cross-section of adopting economies rather than within a single jurisdiction.

We think the paper belongs in this journal specifically because it sits on a question the journal has been building toward from the country side. \textit{Borsa Istanbul Review} has published evidence that green credit policy raises sustainable production across the BRICS economies and that governance quality strengthens that relationship (Jia, Farooq, Alomair, and Aljughaiman, 2026), that China's green-financial-system guidelines drive firms' green transition most where state environmental governance is strong (Chi and Yang, 2023), that green finance improves environmental outcomes only where control of corruption is high (Shi, Zhu, and Khan, 2024), and that regulatory quality conditions rather than merely shifts the financial-development-to-growth relationship (Ullah, Zubir, and Ariff, 2024). Each of those results implies a firm-level prediction: the credit-to-green channel should respond to green banking policy more sharply where regulatory capacity is stronger. Our contribution is to test that prediction directly, on firm-level data, and to report that it does not hold.

Using two independently assembled World Bank Enterprise Survey samples --- a 41-economy harmonized sample on a seven-item green-practice outcome, and a second 162-economy sample reaching into high-income economies entirely outside the network's remit --- we find that bank credit access is robustly associated with green practice adoption, but that neither SBFN adoption nor regulatory quality moderates that relationship under any of four estimators, on either sample. Both institutional variables instead shift the general \emph{level} of green adoption. A country-level Callaway-Sant'Anna staggered event-study on an independent World Development Indicators panel (217 economies, 2000--2024) converges on the same conclusion from a different data source, unit of analysis, and identification strategy. Where effect heterogeneity does exist, the causal forest locates it at firm size rather than country institutions --- which we report as an auxiliary, exploratory finding, and which relocates rather than resolves the institutional question.

We want to disclose one thing directly rather than let a reviewer discover it. An earlier version of this manuscript was submitted elsewhere and desk-rejected, correctly, for two errors in how it described its own identification. First, it asserted that a country fixed effect would absorb the credit-by-SBFN interaction and leave it unidentifiable; that is false, since the interaction is a firm-level regressor and credit access varies within country. Second, it described the Bayesian hierarchical model as the primary test of the paper's three-way hypothesis (H3) when that model contains only two-way cross-level interactions. We have corrected both. No estimate, table, or figure changed: the errors were in prose describing the econometrics, not in the analysis, and H3 was in fact already tested by the saturated triple interaction we now report by name. Section 4.3 now states precisely what a country fixed effect does and does not absorb, and why the primary sample's eight adopter economies make the fixed-effects route imprecise rather than infeasible; Sections 4.3 and 5.3 now scope each estimator to the hypothesis it actually tests. We would rather surface this than have it found.

On identification more broadly, we are explicit in the manuscript (Section 4.1) that the firm-level design is cross-sectional and cannot exploit SBFN's staggered adoption timing, because the Enterprise Surveys' Green Economy module has been fielded to a given country essentially once and never re-fielded comparably after an adoption date. Rather than leave that gap, we built the staggered-adoption design on an independent macro data source and report it in Section 6, including its one significant result and our reasons for not reading it causally. Five tests, three data sources, three units of analysis, and one genuinely quasi-experimental design converge on the same substantive answer.

The manuscript is prepared for double anonymized review: author details appear only in the separate title page file, and the replication repository link is withheld from the manuscript and will be supplied to the editorial office on request. This manuscript is not under consideration elsewhere, and the authors declare no competing interests.

\closing{Sincerely,}

\end{letter}
\end{document}
